\chapter*{Preface}
\addcontentsline{toc}{chapter}{Preface}
\markboth{Preface}{Preface}

Scientific Machine Learning (SciML) is the discipline that fuses the
\emph{inductive bias of physics} with the \emph{expressive flexibility
of modern deep learning}. Where classical numerical analysis treats the
governing equations of nature as the primary object and data as a
secondary product, and where pure machine learning treats data as
primary and equations as optional, SciML insists on the dialogue
between the two. The result is a family of methods --- physics-informed
neural networks, neural operators, Koopman embeddings, universal
differential equations, equivariant networks, diffusion-based
surrogates, foundation models for physics --- that are reshaping how
scientists model complex dynamical systems.

\section*{Who this is for}
This document is written for advanced graduate students and researchers
working at the intersection of dynamical systems, numerical analysis,
and deep learning. It assumes familiarity with multivariable calculus,
linear algebra, ODE/PDE theory, and the basic apparatus of supervised
deep learning (backpropagation, optimisation, common architectures).

\section*{How to read it}
The book is structured as a reference. Part~I (Chapters~1--2) lays the
mathematical foundations of dynamical systems and classical
numerics. Part~II (Chapters~3--7) covers the core SciML toolkit. Part~III
(Chapters~8--11) addresses geometric, generative, and probabilistic
extensions. Part~IV (Chapters~12--13) surveys software, benchmarks,
and emerging trends through the 2024--2026 horizon. The appendix
(Chapter~14) collects functional analysis, automatic differentiation,
and stochastic calculus prerequisites.

Theorems, definitions, recipes, pitfalls, and insights are colour-coded
throughout to help the reader scan rapidly:

\begin{recipe}{Pseudocode and step-by-step procedures}
Practical workflows live in green recipe boxes.
\end{recipe}
\begin{insight}{Conceptual remarks}
Purple boxes flag conceptual subtleties or unifying perspectives.
\end{insight}
\begin{pitfall}{Things that go wrong}
Red boxes highlight failure modes that have caught practitioners.
\end{pitfall}

\section*{Notation}
We follow standard conventions: bold lowercase for vectors
($\bm{x}\in\R^d$), bold uppercase for matrices ($\bm{A}$), calligraphic
for operators and function spaces ($\mathcal{L},\mathcal{G},\mathcal{H}$).
Time is denoted $t$, spatial coordinates $\bm{x}$ or $\bm{\xi}$.
Neural network parameters are $\theta$. Probability density functions
are $p(\cdot)$. Expectations are $\E[\cdot]$. The differential
operator is $\partial$, the gradient $\nabla$, the Laplacian $\Delta$,
and the divergence $\nabla\cdot$.

\section*{Caveats}
SciML moves fast. Where possible we highlight \emph{stable principles}
rather than transient implementations. Citations point to canonical
references; the field's pace means new results may already supersede
specific benchmarks by the time of reading.
